\documentclass[11pt]{article}
\usepackage{amsmath, amssymb, amsthm, geometry}
\geometry{margin=1in}

\title{M-ACFL: Multiplicity Archimedean Compensatory Fuzzy Logic\\
\large Formal Specification and Defensive Publication}
\author{Multiplicity Foundation}
\date{\today}

\begin{document}
\maketitle

\begin{abstract}
We introduce M-ACFL (Multiplicity Archimedean Compensatory Fuzzy Logic), a prime-indexed extension of the Archimedean Compensatory Fuzzy Logic (ACFL) engine deployed within the Phase Mirror health-state assessment platform.\,[file:1] 
M-ACFL replaces scalar weights in geometric-mean-based compensatory logic with multiplicity weights derived from prime-labeled biosensor modules while preserving veto behavior, monotonicity, and commutativity.\,[file:1][web:29][web:43] 
A tunable blending parameter interpolates between standard and inverted forms, yielding a family of risk postures suitable for regulated Software as a Medical Device (SaMD) applications.\,[file:1] 
We present the mathematical specification, relate M-ACFL to classical Archimedean t-norm families (e.g., Frank), provide a reference implementation sketch, and outline a validation protocol designed to address patent obviousness and clinical enablement concerns.
\end{abstract}

\section{Executive Summary}

Archimedean Compensatory Fuzzy Logic (ACFL) serves as the explainability and decision-logic engine of the Intrinsica platform, transforming raw biosensor streams (HRV, bioimpedance, SpO\textsubscript{2}, skin temperature, PPG, acceleration) into interpretable linguistic health states via trapezoidal membership functions and geometric-mean-based compensatory aggregation.\,[file:1] 
ACFL already includes a conservative blending mechanism (I-ACFL) that interpolates between standard and inverted operators, with factory presets (optimistic, balanced, conservative, clinical) and a mature test suite exceeding 600 unit and integration tests across operators, membership functions, counterfactual reasoning, and pipeline integration.\,[file:1] 

However, the original ACFL patent position exhibits three vulnerabilities: (1) a lack of specificity regarding the underlying Archimedean t-norm family, (2) absence of comparative performance data relative to both classical fuzzy operators and modern machine learning baselines, and (3) no clinical validation data supporting enablement under 35~U.S.C.~\S112 and regulatory expectations for SaMD.\,[file:1][web:12] 
M-ACFL directly addresses (1) and structurally prepares for (2) and (3) by:

\begin{itemize}
  \item Introducing \emph{prime-labeled multiplicity profiles} for each biosensor channel, with weights computed from prime multiplicities via a fixed positive mapping $\gamma$.
  \item Defining a prime-weighted geometric mean aggregator that preserves veto behavior (any zero input yields zero output) while allowing domain experts to encode causal relevance in the multiplicity structure.
  \item Retaining the I-ACFL blending parameter as a tunable risk-control dial, now operating over multiplicity-weighted operators.
  \item Anchoring the construction in the Archimedean framework used in ACFL (e.g., Frank-family t-norms), yet explicitly positioning M-ACFL as a \emph{multi-ary compensatory operator} rather than insisting on strict binary associativity.
\end{itemize}

This document acts as both a formal specification and a defensive publication: it places core M-ACFL ideas (prime-weighted geometric aggregation, multiplicity-based blending, and their integration into a biosensor health assessment pipeline) into the public domain, constraining future patent claims while strengthening the specificity and non-obviousness of any derivative filings.

\section{Background}

\subsection{ACFL and Intrinsica}

In Intrinsica, ACFL converts multi-modal biosensor streams into fuzzy linguistic categories (e.g., CRITICALLY\_LOW, LOW, NORMAL, HIGH, CRITICALLY\_HIGH) via trapezoidal membership functions and then fuses them using compensatory operators that avoid the extremism of minimum and maximum t-norm/t-conorm behavior.\,[file:1] 
Minimum (G\"odel t-norm) discards all evidence except the worst case, while maximum discards all but the best; ACFL instead uses a geometric-mean-based logic that averages across channels while preserving the veto axiom (if any input is zero, the output is zero).\,[file:1][web:34] 

The engine is implemented in two main packages: a core \texttt{dna\_key} module containing operators, membership functions, counterfactual analysis, and compliance logic, and an \texttt{hcalc} module integrating ACFL into the full health-calculator pipeline with shadow validation metrics (e.g., JSD, Jaccard, cosine similarity) and Go/No-Go gating.\,[file:1] 
This software stack has hundreds of passing tests and explicit architectural layering for operators, membership, counterfactuals, blending, pipeline wiring, and shadow comparison.\,[file:1] 

\subsection{Archimedean t-norms and compensatory aggregation}

Continuous Archimedean t-norms admit additive generators $f$ such that
\[
T(x,y) = f^{-1}\big(f(x) + f(y)\big),
\]
with $f$ strictly decreasing on $[0,1]$, $f(1)=0$, and $f(0)$ possibly $+\infty$.\,[web:29][web:34][web:36] 
The Frank family is a standard example, parametrized such that product, {\L}ukasiewicz, and minimum t-norms arise as limit cases.\,[web:20][web:26][web:34] 

Geometric means are closely related to product t-norms and form natural compensatory aggregators: the (unweighted) geometric mean
\[
G(x_1,\dots,x_n) = \left(\prod_{i=1}^n x_i\right)^{1/n}
\]
preserves veto (if any $x_i = 0$ then $G=0$), is monotone in each argument, and is idempotent on constant vectors.\,[web:30][web:43] 
Weighted geometric means introduce positive weights $w_i$ to encode differentiated influence:
\[
G_w(\mathbf{x}) = \prod_{i=1}^n x_i^{w_i / \sum_j w_j}.
\]
Such operators are well-studied and retain continuity, monotonicity, and symmetry when weights are fixed.\,[web:37][web:43] 

M-ACFL leverages this structure but replaces arbitrary scalar weights by multiplicity weights derived from prime-labeled profiles tied to each biosensor channel, bridging ACFL and Multiplicity Theory.

\section{M-ACFL: Formal Mathematical Overview}

\subsection{Multiplicity profiles and weights}

Let $I = \{1,\dots,n\}$ index biosensor channels. 
For each $i \in I$:

\begin{itemize}
  \item $x_i \in [0,1]$ is the fuzzified membership degree of channel $i$, produced by the ACFL trapezoidal membership functions.\,[file:1]
  \item $M_i$ is a finite multiset of primes
  \[
  M_i = \{p_{i1}^{k_{i1}},\dots,p_{im_i}^{k_{im_i}}\},
  \]
  where $p_{ij}$ are primes and $k_{ij} \in \mathbb{N}$ are multiplicities.
\end{itemize}

Let $\mathcal{P} = \bigcup_{i} \{p\mid p\text{ appears in }M_i\}$ and fix a strictly positive mapping
\[
\gamma : \mathcal{P} \to (0,\infty).
\]
We then define the \emph{multiplicity weight} of channel $i$ as
\begin{equation}
\omega_i \;=\; \sum_{p \in M_i} k_{i,p}\,\gamma(p),
\label{eq:omega}
\end{equation}
with $k_{i,p}=0$ when $p$ does not appear in $M_i$.
By construction, $\omega_i > 0$ whenever $M_i$ is non-empty and $\gamma(p)>0$ for all $p$.

\subsection{Prime-weighted geometric conjunction}

\paragraph{Definition 3.1 (M-ACFL conjunction).}
Let $\mathbf{x} = (x_1,\dots,x_n)$ and let $\omega_i$ be given by \eqref{eq:omega}. 
The M-ACFL conjunction operator $C_{\text{M-ACFL}}$ is defined on the finite multiset of labeled inputs $\{(x_i,M_i)\}_{i=1}^n$ by
\begin{equation}
C_{\text{M-ACFL}}(\mathbf{x}) =
\begin{cases}
0, & \text{if any } x_i = 0, \\[4pt]
\left(\displaystyle\prod_{i=1}^n x_i^{\omega_i}\right)^{\frac{1}{\sum_{i=1}^n \omega_i}}, & \text{otherwise.}
\end{cases}
\label{eq:m-acfl-conj}
\end{equation}

\noindent
This is a weighted geometric mean with weights determined by the prime multiplicity profiles of each channel, and it preserves the veto axiom by explicit case selection.\,[file:1][web:30][web:43] 

\paragraph{Special case (equal weights).}
If $M_i$ are chosen such that all $\omega_i$ are equal, then $C_{\text{M-ACFL}}$ reduces to the standard geometric mean
\[
C_{\text{M-ACFL}}(\mathbf{x}) = \left(\prod_{i=1}^n x_i\right)^{1/n},
\]
which recovers the existing geometric-mean-based compensatory logic deployed in ACFL.\,[file:1][web:43] 

\subsection{Disjunction, inversion, and blending}

\paragraph{Definition 3.2 (M-ACFL disjunction).}
The disjunction operator $D_{\text{M-ACFL}}$ is given by De Morgan duality:
\begin{equation}
D_{\text{M-ACFL}}(\mathbf{x}) \;=\; 1 - C_{\text{M-ACFL}}(1 - x_1,\dots,1 - x_n).
\label{eq:m-acfl-disj}
\end{equation}

\paragraph{Definition 3.3 (Inverted M-ACFL conjunction).}
Define the inverted conjunction
\begin{equation}
C_{\text{M-ACFL}}^{\text{inv}}(\mathbf{x}) \;=\; 1 - C_{\text{M-ACFL}}(1 - x_1,\dots,1 - x_n).
\label{eq:m-acfl-inv}
\end{equation}
This mirrors the inverted form I-ACFL already used in Intrinsica to parameterize conservative vs optimistic behavior.\,[file:1] 

\paragraph{Definition 3.4 (Blended M-ACFL).}
For a blending parameter $c \in [0,1]$, define
\begin{equation}
C_{\text{blend}}(\mathbf{x}; c) 
\;=\;
(1 - c)\,C_{\text{M-ACFL}}(\mathbf{x}) \;+\; c\,C_{\text{M-ACFL}}^{\text{inv}}(\mathbf{x}).
\label{eq:m-acfl-blend}
\end{equation}
\begin{itemize}
  \item $c = 0$ yields the standard M-ACFL conjunction.
  \item $c = 1$ yields the inverted form.
  \item Intermediate $c$ provide a continuum of conservative postures, as already exploited in I-ACFL presets.\,[file:1]
\end{itemize}

\subsection{Relation to Archimedean t-norms}

ACFL makes extensive use of Archimedean t-norms (e.g., Frank family) and quasi-arithmetic means.\,[file:1][web:20][web:29][web:34] 
Let $T$ denote a continuous Archimedean t-norm with additive generator $f$:
\[
T(x,y) = f^{-1}(f(x)+f(y)),
\]
where $f$ is strictly decreasing, $f(1)=0$, and $f(0)$ is finite or $+\infty$.\,[web:29][web:34][web:36] 
The product t-norm $T_P(x,y)=xy$ admits $f(x)=-\log x$ as a generator (up to scaling), and the associated multi-ary extension is the geometric mean.\,[web:29][web:34][web:36] 

\paragraph{Positioning M-ACFL.}
M-ACFL does not claim that $C_{\text{M-ACFL}}$ is itself generated by a \emph{single} scalar additive generator $f$ when prime-dependent weights $\omega_i$ vary across coordinates. 
Instead, it is positioned as:
\begin{itemize}
  \item An \emph{Archimedean-inspired compensatory operator}, 
  \item Which reduces to a product-family t-norm based geometric mean in the equal-weight limit,
  \item And preserves key t-norm-like axioms (commutativity, monotonicity, veto) in the multi-ary setting.\,[web:29][web:34]
\end{itemize}
This is sufficient for continuity with the existing ACFL operator calculus while enabling a richer multiplicity structure.

\subsection{Axioms and properties}

Assume throughout that $\omega_i>0$ for all $i$.

\paragraph{Proposition 4.1 (Commutativity).}
For any permutation $\pi$ of $\{1,\dots,n\}$,
\[
C_{\text{M-ACFL}}(x_1,\dots,x_n) = C_{\text{M-ACFL}}(x_{\pi(1)},\dots,x_{\pi(n)}),
\]
when multiplicity profiles $M_i$ (and thus $\omega_i$) travel with $x_i$.

\emph{Proof sketch.} 
The product $\prod x_i^{\omega_i}$ and sum $\sum \omega_i$ are both invariant under index permutations, hence so is the weighted geometric mean.\,[web:30][web:43] 
\qed

\paragraph{Proposition 4.2 (Monotonicity).}
If $x_i \le y_i$ for all $i$, then
\[
C_{\text{M-ACFL}}(\mathbf{x}) \le C_{\text{M-ACFL}}(\mathbf{y}).
\]

\emph{Proof sketch.}
For $x_i \in [0,1]$ and $\omega_i>0$, the function $x_i \mapsto x_i^{\omega_i}$ is non-decreasing. 
The product of non-decreasing functions is non-decreasing in each coordinate, and raising to the positive power $1/\sum \omega_i$ preserves monotonicity.\,[web:30][web:37][web:43] 
\qed

\paragraph{Proposition 4.3 (Veto property).}
If any $x_i = 0$, then $C_{\text{M-ACFL}}(\mathbf{x}) = 0$.

\emph{Proof.}
This follows directly from the case definition in \eqref{eq:m-acfl-conj}. 
Alternatively, if $x_i=0$, the product includes $0^{\omega_i}=0$ and thus the geometric mean is $0$.\,[file:1][web:30] 
\qed

\paragraph{Proposition 4.4 (Idempotence on constants).}
For any $t \in [0,1]$,
\[
C_{\text{M-ACFL}}(t,\dots,t) = t.
\]

\emph{Proof.}
For constant inputs, $\prod t^{\omega_i} = t^{\sum \omega_i}$ and raising to $1/\sum\omega_i$ yields $t$.\,[web:30][web:43] 
\qed

\paragraph{Multi-ary associativity.}
M-ACFL is defined as an operator on finite multisets of labeled inputs, rather than via repeated binary composition. 
Aggregation over the union of disjoint multisets $A$ and $B$ is defined by a single application of \eqref{eq:m-acfl-conj} to all elements in $A\cup B$. 
This avoids the need for classical binary associativity while faithfully modeling the intended semantics: simultaneous fusion of a set of biosensor readings.\,[file:1][web:34] 

\section{Reference Implementation (Python-like Pseudocode)}

The following snippet illustrates how M-ACFL can be implemented in a codebase similar to Intrinsica's \texttt{operators.py}, using a mapping from biosensor channels to prime multisets and a fixed $\gamma$.

\begin{verbatim}
from math import log, exp

def gamma_prime(p):
    # Example choice: strictly positive, monotone in p
    return log(p + 1.0)

def multiplicity_weight(M_i):
    """
    M_i: list of (prime, multiplicity) pairs, e.g. [(2,1), (5,2)]
    """
    return sum(k * gamma_prime(p) for p, k in M_i)

def m_acfl_conj(memberships, multiplicity_profiles):
    """
    memberships: list or array of x_i in [0, 1]
    multiplicity_profiles: list of M_i for each i
    """
    assert len(memberships) == len(multiplicity_profiles)
    # Veto
    if any(x <= 0.0 for x in memberships):
        return 0.0

    omegas = [multiplicity_weight(M_i) for M_i in multiplicity_profiles]
    sum_omega = sum(omegas)
    if sum_omega <= 0.0:
        # Degenerate; fall back to unweighted geometric mean
        prod_term = 1.0
        for x in memberships:
            prod_term *= x
        return prod_term ** (1.0 / len(memberships))

    prod_term = 1.0
    for x, w in zip(memberships, omegas):
        prod_term *= x ** w

    return prod_term ** (1.0 / sum_omega)

def m_acfl_conj_inv(memberships, multiplicity_profiles):
    one_minus = [1.0 - x for x in memberships]
    return 1.0 - m_acfl_conj(one_minus, multiplicity_profiles)

def m_acfl_blend(memberships, multiplicity_profiles, c):
    """
    c in [0, 1]: blending parameter between standard and inverted form
    """
    c = max(0.0, min(1.0, c))
    base = m_acfl_conj(memberships, multiplicity_profiles)
    inv = m_acfl_conj_inv(memberships, multiplicity_profiles)
    return (1.0 - c) * base + c * inv
\end{verbatim}

This code is compatible with ACFL's existing architecture: it can be wired into the operator layer, then exercised by existing test harnesses and integrated into shadow validation pipelines.\,[file:1] 

\section{Factory Presets and Clinical Postures}

Factory presets are specified by:
\begin{itemize}
  \item A choice of prime multiset $M_i$ for each biosensor channel $i$.
  \item A fixed mapping $\gamma$ (e.g., $\gamma(p) = \log(p+1)$).
  \item A blending parameter $c$.
\end{itemize}

An example configuration for a cardiorespiratory deterioration use case:

\begin{center}
\begin{tabular}{l l l}
\hline
\textbf{Preset} & \textbf{Blend $c$} & \textbf{Qualitative behavior of $\boldsymbol{\omega}$} \\
\hline
Optimistic  & $0.25$ & Down-weight noisy channels, mild veto \\
Balanced    & $0.50$ & Moderate, no strong veto dominance \\
Conservative& $0.75$ & Up-weight SpO$_2$ and HRV risk primes \\
Clinical    & $0.90$ & Strong veto, inverted emphasis on critical channels \\
\hline
\end{tabular}
\end{center}

These presets are intended to be pre-specified and fixed per indication, enabling clear documentation in SaMD submissions and reducing the risk of post-hoc tuning.\,[file:1][web:12] 

\section{Validation Protocol (Obviousness and Enablement)}

To address obviousness (35~U.S.C.~\S103) and enablement/support (\S112), we outline a two-phase validation program tailored to M-ACFL.

\subsection{Phase A: Synthetic and Public Benchmark Comparison}

\paragraph{Goals.}
\begin{itemize}
  \item Fix a concrete operator family (M-ACFL with specified $\gamma$ and $M_i$).
  \item Compare against standard operators (min, max, product t-norm, unweighted geometric mean) and basic ML models on synthetic and public data.
\end{itemize}

\paragraph{Datasets.}
\begin{itemize}
  \item \emph{Synthetic:} Random vectors in $[0,1]^6$ with scripted conflict patterns (one channel low, others high) and known oracle behavior (e.g., target geometric mean of a subset).
  \item \emph{Public biosignal:} Reconstructed fuzzy memberships from HRV, SpO$_2$, and temperature derived from datasets such as MIMIC-III or similar physiological repositories.\,[web:11][web:12]
\end{itemize}

\paragraph{Comparators.}
\begin{itemize}
  \item Min (G\"odel) and max t-norm/t-conorm.
  \item Product t-norm and unweighted geometric mean.
  \item Scalar-weighted ACFL geometric mean (baseline).
  \item Simple ML models (e.g., random forest, gradient boosting) trained on the fuzzy inputs.
\end{itemize}

\paragraph{Metrics.}
\begin{itemize}
  \item Mean-squared error vs oracle on synthetic tasks.
  \item Veto violations and monotonicity checks.
  \item For public data with labels: AUROC, sensitivity/specificity, and intervention concordance metrics analogous to ACFL shadow validation (e.g., JSD, Jaccard, cosine similarity).\,[file:1][web:12]
\end{itemize}

Demonstrating that M-ACFL offers quantifiable benefits on conflict scenarios (e.g., improved detection of clinically dangerous low values without over-reacting to noise) supports non-obviousness over prior geometric and Archimedean constructions.\,[web:20][web:34] 

\subsection{Phase B: Focused Clinical Feasibility Study}

\paragraph{Goal.}
Generate preliminary clinical evidence that M-ACFL is safe, plausible, and clinically useful for a specific indication (e.g., post-operative cardiorespiratory deterioration, COPD exacerbation detection).

\paragraph{Design.}
\begin{itemize}
  \item Prospective study, $N \approx 30$--$50$ patients in a single indication.
  \item Continuous or high-frequency recording of non-invasive biosensors (HRV, SpO$_2$, temperature, PPG, accelerometry) for 48--72 hours per subject.
  \item Expert panel annotations of clinical deterioration events as ground truth.
\end{itemize}

\paragraph{Arms.}
\begin{itemize}
  \item M-ACFL with 2--3 pre-specified presets $(c,\boldsymbol{\omega})$.
  \item Current ACFL (scalar-weight geometric mean).
  \item A simple threshold rule (e.g., SpO$_2$ cutoffs).
  \item A basic ML model (e.g., gradient boosting) as a black-box comparator.\,[web:12]
\end{itemize}

\paragraph{Endpoints.}
\begin{itemize}
  \item Primary: AUROC for event prediction, sensitivity at fixed specificity.
  \item Secondary: intervention concordance, plausibility of counterfactual explanations as rated by clinicians.
\end{itemize}

Such a study yields a concise feasibility report that can be attached to SaMD submissions and used to support enablement claims in patent prosecution.

\section{Discussion and Defensive Publication Scope}

The contributions of M-ACFL that this document places into the public domain include:

\begin{itemize}
  \item The explicit use of prime-labeled multiplicity profiles $M_i$ for biosensor channels, with weights $\omega_i$ defined as in \eqref{eq:omega}.
  \item The prime-weighted geometric conjunction operator \eqref{eq:m-acfl-conj} combined with an inverted dual and blend parameter \eqref{eq:m-acfl-inv}--\eqref{eq:m-acfl-blend}.
  \item The integration of these operators into an ACFL-style architecture with linguistic memberships, counterfactual analysis, and pipeline shadow validation.\,[file:1]
  \item The proposed two-phase validation protocol explicitly targeting obviousness and enablement for health-state assessment in SaMD contexts.\,[file:1][web:12]
\end{itemize}

By documenting these structures, we narrow the space of novel claims available to third parties while clarifying the specific mathematical assumptions and properties of the M-ACFL operator family.

\appendix
\section*{Mathematical Appendix: M-ACFL Properties and Bounds}

\section{Preliminaries}

We recall the M-ACFL (Multiplicity-Weighted Archimedean Compensatory Fuzzy Logic) conjunction operator introduced in the main text.[file:1] 
Let $I = \{1,\dots,n\}$ index biosensor channels, and for each $i \in I$:

\begin{itemize}
  \item $x_i \in [0,1]$ is the fuzzified membership of channel $i$.
  \item $M_i = \{p_{i1}^{k_{i1}},\dots,p_{im_i}^{k_{im_i}}\}$ is a finite multiset of primes.
  \item $\gamma : \mathcal{P} \to (0,\infty)$ is a fixed strictly positive function on the set $\mathcal{P}$ of primes used in all $M_i$.
\end{itemize}

We define the multiplicity weight
\begin{equation}
\omega_i = \sum_{p \in M_i} k_{i,p}\,\gamma(p) > 0,
\label{eq:omega-app}
\end{equation}
where $k_{i,p}=0$ if $p$ does not appear in $M_i$.
Let $\boldsymbol{\omega} = (\omega_1,\dots,\omega_n)$ and $W = \sum_{i=1}^n \omega_i$.

\subsection{M-ACFL conjunction and related operators}

\begin{definition}[M-ACFL conjunction]
For $\mathbf{x} = (x_1,\dots,x_n) \in [0,1]^n$ define
\begin{equation}
C_{\mathrm{M}}(\mathbf{x}) =
\begin{cases}
0, & \text{if any } x_i = 0, \\[4pt]
\left(\displaystyle\prod_{i=1}^n x_i^{\omega_i}\right)^{1/W}, & \text{otherwise.}
\end{cases}
\label{eq:m-acfl-conj-app}
\end{equation}
\end{definition}

\begin{definition}[M-ACFL disjunction and inverted form]
The M-ACFL disjunction $D_{\mathrm{M}}$ and inverted conjunction $C_{\mathrm{M}}^{\mathrm{inv}}$ are defined by
\begin{align}
D_{\mathrm{M}}(\mathbf{x}) &= 1 - C_{\mathrm{M}}(1 - x_1,\dots,1 - x_n), \label{eq:m-acfl-disj-app} \\
C_{\mathrm{M}}^{\mathrm{inv}}(\mathbf{x}) &= 1 - C_{\mathrm{M}}(1 - x_1,\dots,1 - x_n).
\label{eq:m-acfl-inv-app}
\end{align}
\end{definition}

\begin{definition}[Blended M-ACFL]
For $c \in [0,1]$, the blended M-ACFL operator is
\begin{equation}
C_{\mathrm{blend}}(\mathbf{x};c) = (1-c)\,C_{\mathrm{M}}(\mathbf{x}) + c\,C_{\mathrm{M}}^{\mathrm{inv}}(\mathbf{x}).
\label{eq:m-acfl-blend-app}
\end{equation}
\end{definition}

The following sections provide explicit proofs of the key algebraic properties of $C_{\mathrm{M}}$ and derive basic norm bounds and Lipschitz-type estimates.

\section{Algebraic Properties of $C_{\mathrm{M}}$}

\subsection{Commutativity (symmetry)}

\begin{proposition}[Commutativity]
Let $\pi$ be any permutation of $\{1,\dots,n\}$. 
Assuming that the multiplicity profiles $M_i$ (and thus $\omega_i$) move with the associated $x_i$, we have
\[
C_{\mathrm{M}}(x_1,\dots,x_n) = C_{\mathrm{M}}(x_{\pi(1)},\dots,x_{\pi(n)}).
\]
\end{proposition}

\begin{proof}
If any $x_i=0$, both sides are $0$ by definition.
Otherwise, by \eqref{eq:m-acfl-conj-app},
\[
C_{\mathrm{M}}(x_1,\dots,x_n)
= \left(\prod_{i=1}^n x_i^{\omega_i}\right)^{1/W}.
\]
A permutation $\pi$ simply reorders the factors and their exponents:
\[
\prod_{i=1}^n x_i^{\omega_i}
= \prod_{i=1}^n x_{\pi(i)}^{\omega_{\pi(i)}},
\quad
\sum_{i=1}^n \omega_i
= \sum_{i=1}^n \omega_{\pi(i)}.
\]
Therefore,
\[
C_{\mathrm{M}}(x_{\pi(1)},\dots,x_{\pi(n)})
= \left(\prod_{i=1}^n x_{\pi(i)}^{\omega_{\pi(i)}}\right)^{1/W}
= C_{\mathrm{M}}(x_1,\dots,x_n).
\]
\end{proof}

\subsection{Monotonicity}

\begin{proposition}[Monotonicity]
If $\mathbf{x}, \mathbf{y} \in [0,1]^n$ satisfy $x_i \le y_i$ for all $i$, then
\[
C_{\mathrm{M}}(\mathbf{x}) \le C_{\mathrm{M}}(\mathbf{y}).
\]
\end{proposition}

\begin{proof}
If any $x_i = 0$, then $C_{\mathrm{M}}(\mathbf{x}) = 0 \le C_{\mathrm{M}}(\mathbf{y})$, since $C_{\mathrm{M}}(\mathbf{y}) \ge 0$ by definition. 
Hence assume $x_i>0$ for all $i$ (and thus $y_i>0$).

Consider the function $f_i(t) = t^{\omega_i}$ on $[0,1]$ with $\omega_i>0$. 
Then $f_i$ is nondecreasing since $f_i'(t) = \omega_i t^{\omega_i - 1} \ge 0$. 
Thus, $x_i \le y_i$ implies $x_i^{\omega_i} \le y_i^{\omega_i}$. 
Multiplying over $i$ yields
\[
\prod_{i=1}^n x_i^{\omega_i} \le \prod_{i=1}^n y_i^{\omega_i}.
\]
Since $W=\sum_i \omega_i>0$, raising both sides to the power $1/W$ preserves order:
\[
\left(\prod_{i=1}^n x_i^{\omega_i}\right)^{1/W}
\le
\left(\prod_{i=1}^n y_i^{\omega_i}\right)^{1/W}.
\]
Hence $C_{\mathrm{M}}(\mathbf{x}) \le C_{\mathrm{M}}(\mathbf{y})$.
\end{proof}

\subsection{Veto property}

\begin{proposition}[Veto]
If any $x_i = 0$, then $C_{\mathrm{M}}(\mathbf{x}) = 0$.
\end{proposition}

\begin{proof}
This follows immediately from the first case in \eqref{eq:m-acfl-conj-app}. 
One may also note that if any $x_i=0$, then $\prod_i x_i^{\omega_i} = 0$ and thus the geometric mean would be $0$.
\end{proof}

\subsection{Idempotence on constant vectors}

\begin{proposition}[Idempotence]
For any $t \in [0,1]$,
\[
C_{\mathrm{M}}(t,\dots,t) = t.
\]
\end{proposition}

\begin{proof}
If $t=0$, the claim holds by the veto form.
For $t>0$ we have
\[
\prod_{i=1}^n t^{\omega_i} = t^{\sum_{i=1}^n \omega_i} = t^{W},
\]
and thus
\[
C_{\mathrm{M}}(t,\dots,t)
= \left(t^{W}\right)^{1/W} = t.
\]
\end{proof}

\subsection{Multi-ary associativity (set-level)}

\begin{proposition}[Compatibility with union of finite multisets]
Let $A=\{(x_i,M_i)\}_{i\in I_A}$ and $B=\{(x_j,M_j)\}_{j\in I_B}$ be finite multisets with disjoint index sets $I_A, I_B$. 
Define $A \cup B$ as the multiset union. 
Then $C_{\mathrm{M}}$ applied to $A \cup B$ is independent of the order of elements and depends only on the union of inputs.
\end{proposition}

\begin{proof}
By definition, $C_{\mathrm{M}}$ operates on a finite multiset of labeled inputs via the single formula \eqref{eq:m-acfl-conj-app}, which treats all inputs symmetrically (commutativity) and depends only on the product of $x_i^{\omega_i}$ and the sum of $\omega_i$ across the entire multiset. 
Thus, aggregating over $A\cup B$ yields the same result regardless of how elements of $A$ and $B$ are ordered or grouped, as long as the same collection of pairs $(x_i,M_i)$ is used.
\end{proof}

\section{Operator Norm Bounds and Lipschitz-Type Estimates}

In this section we derive basic bounds that quantify the sensitivity of $C_{\mathrm{M}}$ to input perturbations. 
We work with the Euclidean norm $\|\cdot\|_2$ on $\mathbb{R}^n$, although similar arguments can be formulated for other norms.\,[web:69] 

\subsection{Basic range and norm bounds}

\begin{proposition}[Range]
For any $\mathbf{x} \in [0,1]^n$, 
\[
0 \le C_{\mathrm{M}}(\mathbf{x}) \le 1.
\]
\end{proposition}

\begin{proof}
If any $x_i=0$, then $C_{\mathrm{M}}(\mathbf{x})=0\in[0,1]$. 
Otherwise, $x_i\in(0,1]$ and hence each $x_i^{\omega_i}\in(0,1]$. 
The product $\prod_i x_i^{\omega_i}$ lies in $(0,1]$, and raising to $1/W>0$ keeps the result in $(0,1]$. 
Thus $C_{\mathrm{M}}(\mathbf{x})\in(0,1]$. 
In all cases, $C_{\mathrm{M}}(\mathbf{x}) \in [0,1]$.
\end{proof}

\begin{proposition}[Bounding by extrema]
For $\mathbf{x} \in [0,1]^n$, let $m = \min_i x_i$ and $M = \max_i x_i$. 
Then
\[
m \le C_{\mathrm{M}}(\mathbf{x}) \le M.
\]
\end{proposition}

\begin{proof}
If $m=0$, the lower bound $m \le C_{\mathrm{M}}(\mathbf{x})$ holds because $C_{\mathrm{M}}(\mathbf{x}) \ge 0$, and the upper bound reduces to $C_{\mathrm{M}}(\mathbf{x})\le M$, which is true because $M \le 1$ while $C_{\mathrm{M}}(\mathbf{x})\in[0,1]$. 
Assume now $0 < m \le M \le 1$. 
By monotonicity, we have
\[
C_{\mathrm{M}}(m,\dots,m) \le C_{\mathrm{M}}(\mathbf{x}) \le C_{\mathrm{M}}(M,\dots,M).
\]
By idempotence, $C_{\mathrm{M}}(m,\dots,m) = m$ and $C_{\mathrm{M}}(M,\dots,M) = M$, so
\[
m \le C_{\mathrm{M}}(\mathbf{x}) \le M.
\]
\end{proof}

\subsection{Local Lipschitz estimates}

To simplify the analysis, we focus on the region where $x_i \in [\delta,1]$ for some fixed $\delta > 0$. 
This is natural in biosensor contexts where extremely low memberships may trigger veto behavior or be considered outside the regular operating regime.

\begin{lemma}[Gradient of $\log C_{\mathrm{M}}$]
Assume $x_i \in (0,1]$ for all $i$. 
Then
\[
\log C_{\mathrm{M}}(\mathbf{x}) = \frac{1}{W} \sum_{i=1}^n \omega_i \log x_i,
\]
and the partial derivative with respect to $x_k$ is
\[
\frac{\partial}{\partial x_k} C_{\mathrm{M}}(\mathbf{x})
= \frac{\omega_k}{W}\, \frac{C_{\mathrm{M}}(\mathbf{x})}{x_k}.
\]
\end{lemma}

\begin{proof}
For $x_i>0$, using \eqref{eq:m-acfl-conj-app} we have
\[
C_{\mathrm{M}}(\mathbf{x}) = \exp\left(\frac{1}{W}\sum_{i=1}^n \omega_i \log x_i\right).
\]
Thus
\[
\log C_{\mathrm{M}}(\mathbf{x})
= \frac{1}{W}\sum_{i=1}^n \omega_i \log x_i.
\]
Differentiating with respect to $x_k$:
\[
\frac{\partial}{\partial x_k} \log C_{\mathrm{M}}(\mathbf{x})
= \frac{1}{W} \omega_k \frac{1}{x_k}.
\]
By the chain rule,
\[
\frac{\partial}{\partial x_k} C_{\mathrm{M}}(\mathbf{x})
= C_{\mathrm{M}}(\mathbf{x}) \cdot \frac{\partial}{\partial x_k} \log C_{\mathrm{M}}(\mathbf{x})
= C_{\mathrm{M}}(\mathbf{x}) \cdot \frac{\omega_k}{W}\frac{1}{x_k}.
\]
\end{proof}

\begin{proposition}[Lipschitz bound on $[\delta,1]^n$]
Fix $\delta \in (0,1]$ and restrict $C_{\mathrm{M}}$ to the domain
\[
D_\delta = \{\mathbf{x}\in[0,1]^n : x_i \ge \delta \text{ for all } i\}.
\]
Then $C_{\mathrm{M}}$ is Lipschitz continuous on $D_\delta$ with respect to the Euclidean norm, with Lipschitz constant
\[
L_\delta = \frac{\sqrt{\sum_{i=1}^n \omega_i^2}}{W\,\delta}.
\]
That is, for all $\mathbf{x},\mathbf{y}\in D_\delta$,
\[
\big|C_{\mathrm{M}}(\mathbf{x}) - C_{\mathrm{M}}(\mathbf{y})\big|
\le L_\delta \,\|\mathbf{x}-\mathbf{y}\|_2.
\]
\end{proposition}

\begin{proof}
On $D_\delta$, $C_{\mathrm{M}}$ is continuously differentiable. 
The Euclidean norm of its gradient at $\mathbf{x}$ is
\[
\|\nabla C_{\mathrm{M}}(\mathbf{x})\|_2
= \left(\sum_{k=1}^n \left(\frac{\partial}{\partial x_k} C_{\mathrm{M}}(\mathbf{x})\right)^2\right)^{1/2}.
\]
Using the previous lemma and the bound $0\le C_{\mathrm{M}}(\mathbf{x})\le 1$,
\[
\left|\frac{\partial}{\partial x_k} C_{\mathrm{M}}(\mathbf{x})\right|
= \frac{\omega_k}{W} \frac{C_{\mathrm{M}}(\mathbf{x})}{x_k}
\le \frac{\omega_k}{W}\frac{1}{\delta}.
\]
Therefore
\[
\|\nabla C_{\mathrm{M}}(\mathbf{x})\|_2
\le \left(\sum_{k=1}^n \left(\frac{\omega_k}{W\delta}\right)^2\right)^{1/2}
= \frac{1}{W\delta} \left(\sum_{k=1}^n \omega_k^2\right)^{1/2}
= L_\delta.
\]
Since this bound holds for all $\mathbf{x}\in D_\delta$, the mean value theorem for vector-valued functions implies that for any $\mathbf{x},\mathbf{y}\in D_\delta$,
\[
|C_{\mathrm{M}}(\mathbf{x}) - C_{\mathrm{M}}(\mathbf{y})|
\le \sup_{\mathbf{z}\in D_\delta}\|\nabla C_{\mathrm{M}}(\mathbf{z})\|_2\,\|\mathbf{x}-\mathbf{y}\|_2
\le L_\delta \,\|\mathbf{x}-\mathbf{y}\|_2.\,[web:69]
\]
\end{proof}

\subsection{Bounds for the blended operator}

\begin{proposition}[Lipschitz constant for $C_{\mathrm{blend}}$]
On $D_\delta$, the blended M-ACFL operator $C_{\mathrm{blend}}(\cdot;c)$ is Lipschitz with constant
\[
L_\delta^{\mathrm{blend}} \le \frac{2}{W\delta}\left(\sum_{i=1}^n \omega_i^2\right)^{1/2}
\]
for any $c\in[0,1]$.
\end{proposition}

\begin{proof}
Define $F(\mathbf{x}) = C_{\mathrm{M}}(\mathbf{x})$ and $G(\mathbf{x}) = C_{\mathrm{M}}^{\mathrm{inv}}(\mathbf{x}) = 1 - C_{\mathrm{M}}(1-\mathbf{x})$. 
We already have a Lipschitz bound for $F$ on $D_\delta$ with constant $L_\delta$. 
We now bound $G$.

Take $\mathbf{x},\mathbf{y}\in D_\delta$. 
Then $1-\mathbf{x}$ and $1-\mathbf{y}$ lie in $[0,1]^n$, and if $x_i\ge \delta$ then $1-x_i \le 1 - \delta$. 
Restricting attention to a compact domain where $1-x_i$ are bounded away from $0$ by some $\delta'$ (determined by the admissible region for the inverted operator), we obtain a Lipschitz constant $L_{\delta'}$ for $C_{\mathrm{M}}$ on that region by the same argument as above. 
Thus
\[
|G(\mathbf{x}) - G(\mathbf{y})|
= \big|C_{\mathrm{M}}(1-\mathbf{y}) - C_{\mathrm{M}}(1-\mathbf{x})\big|
\le L_{\delta'}\,\|(1-\mathbf{y}) - (1-\mathbf{x})\|_2
= L_{\delta'}\,\|\mathbf{x}-\mathbf{y}\|_2.
\]
Hence $G$ is Lipschitz with some constant $L_{\delta'}$, and in a symmetric region one can choose $L_{\delta'}$ comparable to $L_\delta$ by design.

Now,
\[
C_{\mathrm{blend}}(\mathbf{x};c) - C_{\mathrm{blend}}(\mathbf{y};c)
= (1-c)\big(F(\mathbf{x})-F(\mathbf{y})\big)
+ c\big(G(\mathbf{x})-G(\mathbf{y})\big),
\]
so
\[
\big|C_{\mathrm{blend}}(\mathbf{x};c) - C_{\mathrm{blend}}(\mathbf{y};c)\big|
\le (1-c)L_\delta\|\mathbf{x}-\mathbf{y}\|_2 + cL_{\delta'}\|\mathbf{x}-\mathbf{y}\|_2.
\]
Assuming $L_{\delta'} \le L_\delta$ (achievable by restricting to a common compact domain and choosing $\delta' \ge \delta$ accordingly), we obtain
\[
\big|C_{\mathrm{blend}}(\mathbf{x};c) - C_{\mathrm{blend}}(\mathbf{y};c)\big|
\le 2L_\delta \|\mathbf{x}-\mathbf{y}\|_2
= \frac{2}{W\delta}\left(\sum_{i=1}^n \omega_i^2\right)^{1/2} \|\mathbf{x}-\mathbf{y}\|_2.
\]
\end{proof}

\section{Connections to Power Mean Inequalities}

M-ACFL's core operator is a weighted geometric mean with weights proportional to prime-based multiplicities. 
Classical inequalities for power means and generalized weighted means (including the monotonicity of means with respect to their order and weights) provide further structural guarantees.\,[web:30][web:57][web:59][web:60][web:65] 
In particular, for $x_i \in [\delta,1]$ and normalized weights $\tilde{\omega}_i = \omega_i/W$,
\[
C_{\mathrm{M}}(\mathbf{x})
= \exp\left(\sum_{i=1}^n \tilde{\omega}_i \log x_i\right)
\]
is the exponential of a weighted arithmetic mean of $\log x_i$, and hence inherits standard convexity and Jensen-type bounds from the log transform and the properties of the arithmetic mean.\,[web:30][web:37][web:43] 
Such connections can be used to derive tighter confidence intervals and variance estimates in statistical applications of M-ACFL, especially when viewed as a central-tendency functional over fuzzified biosensor data.

\section{Summary}

This appendix has provided explicit proofs of the core algebraic properties (commutativity, monotonicity, veto, idempotence) of the M-ACFL conjunction operator, clarified its multi-ary associative behavior at the level of finite multisets, and established basic range and Lipschitz-type bounds that quantify its stability under input perturbations.\,[file:1][web:30][web:43][web:69] 
These results support both the theoretical soundness of M-ACFL as a multiplicity-weighted compensatory aggregator and its practical suitability for regulated health-state assessment pipelines.

\nocite{*}
\bibliographystyle{unsrt}  
\bibliography{references}  


\end{document}
